\documentclass[10pt,a4paper]{article}
\usepackage[T1]{fontenc}
\usepackage[utf8]{inputenc}
\usepackage{amsmath,amssymb}
\usepackage{geometry}
\usepackage{microtype}
\usepackage{lmodern}
\usepackage[hidelinks]{hyperref}
\geometry{margin=2.15cm}

\setlength{\parindent}{0pt}
\setlength{\parskip}{0.35em}

\hypersetup{
  pdftitle={Cover letter: Exact Affine Geometry, Optimal Centres, Certified Compression Bounds, and Flag-Quotient Width Obstructions for Finite-Outcome Quantum Instruments},
  pdfauthor={Amin Abaee},
  pdfsubject={Cover letter to Quantum Information Processing},
  pdfcreator={LaTeX via Tectonic}
}

\begin{document}
\pagestyle{empty}

\noindent
\begin{minipage}[t]{0.60\textwidth}
Amin Abaee\\
Independent Researcher\\
Email: \href{mailto:amin\_abaee@ut.ac.ir}{amin\_abaee@ut.ac.ir}\\
ORCID: 0000-0002-0019-1842
\end{minipage}%
\hfill
\begin{minipage}[t]{0.36\textwidth}
\raggedleft
September 14, 2026
\end{minipage}

\vspace{1.0em}

\noindent The Editors\\
\emph{Quantum Information Processing}\\
Springer Nature

\vspace{1.0em}

\noindent\textbf{Re:} Submission of the research article \emph{Exact Affine Geometry, Optimal Centres, Certified Compression Bounds, and Flag-Quotient Width Obstructions for Finite-Outcome Quantum Instruments}

\vspace{1.0em}

\noindent Dear Editors,

I am pleased to submit the manuscript cited above --- a 53-page, single-author research article --- for consideration in \emph{Quantum Information Processing}.

The manuscript determines, exactly, the geometry of the body of finite-outcome quantum instruments --- the objects describing a measurement's full statistics together with its conditional post-measurement states, and the building blocks of adaptive protocols and quantum networks --- and derives its compression consequences. The affine dimension of the body of $n$-outcome instruments between systems of dimensions $d_A$ and $d_B$ is $d_A^2(nd_B^2-1)$; the largest diamond-norm ball inside it has radius $2/(nd_B\,\min\{d_A,d_B\})$, is centered at the uniform depolarising instrument, and the same instrument is simultaneously optimal for the covering problem, the two radii linked by a complementarity identity. Every radius and width carries a single-shot operational meaning through the optimal discrimination error. On the compression side, antipodal constructions certify a lower envelope on the optimal error, while constructive upper bounds from constant codes, convex projection, and balanced pinching meet the envelope exactly at zero latent dimension, in the collapse case, and from the affine dimension on, where the zero-error threshold equals the affine dimension. For input-independent bodies the flat-width classification is complete: flat precisely on the classical ($n\le 4$) and qubit ($n\le 2$) bodies, bounded below by $4/3$ on every other, with the qutrit exactly $4/3$ --- the second-index result resting on the integral cohomology of the cyclic flag quotient, cross-checked by an independent exact-integer machine computation whose code and transcripts are publicly available.

Instruments are the operational primitives of adaptive quantum protocols --- measurement and state update in a single object --- so the parameter cost of representing them faithfully is a basic resource question for quantum information. The manuscript is exact where the neighbouring literature is asymptotic: it maintains an explicit distinction between quantities determined exactly and those bounded by certified two-sided envelopes, and its cohomological obstructions --- Borsuk--Ulam arguments and flag-quotient cup products --- reach impossibility results that metric arguments alone do not, with the precise relation to the existing cohomological literature recorded in the text.

The general source--latent theory underlying the obstruction arguments --- fiber radii, antipodal profiles, the deleted-product coindex --- is developed in a companion paper on quantum channels, \emph{Exact Diamond Balls, Antipodal Widths, and Query-Uniform Compression Bounds for Quantum Channels}, which is being submitted separately to the journal; the two papers are self-contained, cite each other, and can be reviewed independently.

The submission lies squarely within the journal's scope of theoretical research across quantum information science. The questions --- how many parameters an instrument body retains under continuous encoding, and which bodies admit flat, error-free encodings --- belong to channel and instrument simulation and to resource-efficient representation; the formalism is standard, the article is full-length and self-contained, and the machine computation certifying one cohomological input is reproducible and publicly hosted.

I confirm that the manuscript is original, unpublished, and not under consideration elsewhere; I am its sole author, with no funding and no competing interests to declare. The supporting code and computation transcripts are available at the public repository cited in the manuscript.

Thank you for your consideration. I look forward to the outcome of the review.

\vspace{0.8em}

\noindent Yours sincerely,

\vspace{1.4em}

\noindent
\begin{tabular}{@{}l@{}}
Amin Abaee\\
Independent Researcher\\
Email: \href{mailto:amin\_abaee@ut.ac.ir}{amin\_abaee@ut.ac.ir}
\end{tabular}

\end{document}
